\section{Literature Review}
\label{ch:lit}

This section summarises the work on which each part of the project rests, in the order in which
the ideas are used.

\subsection{Mean-variance portfolio theory}

\citet{markowitz1952} showed that an investor who cares about expected return and variance should
not judge an asset on its own, but by what it adds to the whole portfolio. Assets that do not
move together reduce risk when combined; this diversification effect is the reason for holding Nifty
and gold together (Section~\ref{sec:maxsharpe}).

\citet{tobin1958} added a risk-free asset and showed that the problem separates into two decisions:
(i) the mix of risky assets, which is the same for every investor, and (ii) how much to hold in the
risk-free asset, which depends on the investor's attitude to risk. This \emph{separation} is the
basis of the Personal SIP (Section~\ref{sec:personal}): the liquid fund plays the role of the
risk-free asset, its share is set by the investor's risk aversion, and the Nifty/gold split is
common to everyone.

\subsection{Risk-adjusted performance: the Sharpe ratio}

\citet{sharpe1966} proposed the reward-to-variability ratio, later called the Sharpe ratio: the
return in excess of a risk-free rate divided by volatility. \citet{sharpe1994} restated it and
stressed that the benchmark rate must be the investor's actual alternative. This point turned out
to be central in this project: with a liquid fund in the portfolio, measuring excess return against
0\% instead of against the liquid fund makes cash look nearly perfect
(Section~\ref{sec:res-2010}).

\subsection{Risk aversion and mean-variance utility}

\citet{pratt1964} formalised \emph{risk aversion} and the \emph{certainty equivalent}: the sure
amount an investor would accept in place of a risky prospect. For returns that are not extreme,
expected utility is closely approximated by a function of mean and variance alone, and
\citet{levy1979} and \citet{kroll1984} showed empirically that portfolios chosen by such a
mean-variance score are almost identical to those chosen by maximising expected utility directly.
The standard textbook form \citep{bkm2021} is
\begin{equation*}
  U = E(r) - \tfrac{1}{2}A\sigma^2 ,
\end{equation*}
where $A$ is the risk-aversion coefficient. This project uses the same form with returns measured
above the liquid fund and $A$ written as $\lambda$ (Section~\ref{sec:utility}). Maximising $U$ is
equivalent to Markowitz's problem of maximising return for a given variance, with $\lambda$ acting
as the Lagrange multiplier on the variance constraint.

\subsection{Risk parity}

Risk parity allocates so that each asset contributes the same share of total portfolio risk.
\citet{maillard2010} define the risk contribution of asset $i$ as $w_i(\Sigma w)_i$ and study the
equal-risk-contribution portfolio; \citet{qian2006} shows that these contributions add up to total
variance and can be read as each asset's share of likely losses. \citet{asness2012} explain why such
portfolios have historically delivered good risk-adjusted returns. Because risk parity uses only
volatilities and correlations, not expected returns, it is far less sensitive to estimation error,
which is why it is one of the two components of the allocation rule used here.

\subsection{Estimation error}

Optimisers need estimates of expected returns, variances and covariances, and these are noisy.
\citet{michaud1989} showed that mean-variance optimisers act as ``error maximisers'': they put the
most weight on assets whose returns happen to be overestimated. \citet{chopra1993} found that errors
in expected returns are roughly ten times as damaging as errors in variances. \citet{demiguel2009}
found that a simple $1/N$ (equal-weight) portfolio was not consistently beaten out of sample by
fourteen optimised models, while \citet{kritzman2010} argued that optimisation does add value when
inputs are estimated sensibly. This debate is tested directly in Section~\ref{ch:results}, where an
equal-weight SIP, several optimisers and a direct utility maximiser are compared on unseen data.
\citet{bailey2014} warn that trying many variations on one data set produces back-tests that look
good by chance; the strict separation of design (2010--2019) and test (2020--2025) periods in this
project follows that advice.

\subsection{Gold, SIPs and rebalancing}

Gold is a hedge against equities on average and a safe haven in market crashes \citep{baur2010}; for an Indian investor, rupee gold also rises when the rupee weakens. 
A SIP is a form of rupee-cost averaging, which is not optimal for a lump sum \citep{constantinides1979,vanduffel2012}; for a salaried investor with no lump sum, the question becomes how to allocate each instalment. How often to rebalance is studied by \citet{perold1988}, who compare constant-mix and other dynamic
strategies, and by \citet{jaconetti2010}, who recommend rebalancing only when allocations drift
beyond a threshold, in the range of about 5\%, rather than on a fixed calendar. The 5\% band rule in
this project's engine follows that recommendation.

